% Generated draft: requires mathematical and source review.
\documentclass[11pt]{article}
\usepackage[margin=25mm]{geometry}
\usepackage{fontspec}
\setmainfont{Latin Modern Roman}
\usepackage{amsmath,amssymb,mathtools,unicode-math}
\setmathfont{Latin Modern Math}
\usepackage{enumitem,needspace,xurl,hyperref}
\setlength{\parindent}{0pt}
\setlength{\parskip}{5pt}
\setlength{\emergencystretch}{3em}
\newcommand{\N}{\mathbb{N}}
\newcommand{\Z}{\mathbb{Z}}
\newcommand{\Q}{\mathbb{Q}}
\newcommand{\R}{\mathbb{R}}
\newcommand{\C}{\mathbb{C}}
\newcommand{\F}{\mathbb{F}}
\newcommand{\A}{\mathbb{A}}
\newcommand{\PP}{\mathbb{P}}
\newcommand{\E}{\mathbb{E}}
\newcommand{\eps}{\varepsilon}
\newcommand{\dd}{\,\mathrm{d}}
\begin{document}
\section*{UnsolvedMath ID: 30002520}
\textbf{Source problem code:} catalogue entry
\subsection{Definitions and mathematical statement}
Let a partition be a finite weakly decreasing sequence of positive integers, with size equal to the sum of its parts. Write \(\mu=(\mu^1,\ldots,\mu^k)\) for a \(k\)-tuple of partitions. For a box \(s\) in the Young diagram of a partition \(\lambda\), let \(a(s)\) and \(l(s)\) denote its arm length and leg length, respectively: the numbers of boxes strictly to its right and strictly below it.

The supplied record associates to \(\mu\) a generic \(\operatorname{GL}_n(\C)\) character variety \(\mathcal M_\mu\). It does not define this variety, the genericity conditions, or the relation between \(n\) and the sizes of the partitions. In particular, it does not explicitly state the usual common-size requirement \(|\mu^1|=\cdots=|\mu^k|=n\). The parameter \(g\) appears in the following formula without an explicit definition or range in the record.

With \(z,w\) formal variables, the supplied generating-function prescription is
\[
\mathcal H_\lambda(z,w)
=\prod_{s\in\lambda}
\frac{\bigl(z^{2a(s)+1}-w^{2l(s)+1}\bigr)^{2g}}
{\bigl(z^{2a(s)+2}-w^{2l(s)}\bigr)
 \bigl(z^{2a(s)}-w^{2l(s)+2}\bigr)},
\]
\[
\begin{aligned}
&(z^2-1)(1-w^2)\operatorname{Log}\!\left(
\sum_\lambda \mathcal H_\lambda(z,w)
\prod_{i=1}^k\widetilde H_\lambda(x_i;z^2,w^2)
\right)\\
&\hspace{3em}=\sum_\mu\mathbb H_\mu(z,w)m_\mu(x).
\end{aligned}
\]
Here the record calls \(\widetilde H_\lambda\) Macdonald polynomials. Under the conventional interpretation of the symmetric-function notation, \(x_1,\ldots,x_k\) are separate alphabets and
\[
m_\mu(x)=\prod_{i=1}^k m_{\mu^i}(x_i),
\]
where \(m_\nu\) is the monomial symmetric function indexed by \(\nu\). This interpretation requires confirmation, as do the normalization of \(\widetilde H_\lambda\), the meaning of \(\operatorname{Log}\), and the formal completion and indexing used in the expansion. In particular, the record does not explicitly specify whether \(\operatorname{Log}\) is the plethystic logarithm.

\textbf{Conjectural target as supplied.} For every tuple \(\mu\) and every choice of parameters satisfying the intended, but incompletely specified, character-variety hypotheses, are the functions \(\mathbb H_\mu(z,w)\) polynomials in \(z,w\) that encode the mixed Hodge numbers of \(\mathcal M_\mu\)?

The polynomiality clause is an explicit target. The mixed-Hodge clause is not a complete mathematical identity in the supplied record: no substitution, grading convention, dimension shift, or choice between ordinary and compactly supported cohomology is given. Consequently, the displayed prescription and the stated conjectural target are preserved here, but a precise universal mixed-Hodge formula cannot be supplied without additional source information. These unresolved qualifications are part of this statement-only draft.
\subsection{Short English statement}
The supplied conjecture asks whether coefficients \(\mathbb H_\mu(z,w)\), obtained from a Macdonald-polynomial generating function, are polynomials describing the mixed Hodge numbers of generic \(\operatorname{GL}_n(\C)\) character varieties. The record omits the exact character-variety hypotheses and the identity relating these coefficients to mixed Hodge numbers, so the full conjecture requires source clarification and mathematical review.
\subsection{Sources}
\begin{itemize}
\item Supplied local UnsolvedMath registry record, permanent ID \(30002520\), title \emph{Mixed Hodge Polynomials of Character Varieties}, exact source problem code \texttt{catalogue entry}, in the collection \emph{Oberwolfach Reports Open Problems}. The record gives proposed year 2014 and names no individual proposer.
\item Supplied report citation: \emph{Representation Theory of Quivers and Finite Dimensional Algebras} (2014), Oberwolfach Reports, pp.\ 453--529, DOI 10.4171/owr/2014/08. \url{https://doi.org/10.4171/owr/2014/08}. The background identifies Oberwolfach Reports as reports of mathematical workshops and their open-problem discussions. No exact problem page or passage is supplied.
\item Supplied literature source: Hausel--Letellier--Rodriguez-Villegas. \url{https://arxiv.org/abs/0810.2076}.
\item Supplied literature source: Mellit, integrality of HLRV kernels. \url{https://www.mellit.org/publication/mellit-2018-integrality/}.
\item Literature assessment dated exactly \texttt{2026-08-21}: the supplied curation labels the problem partially solved and states that E-polynomial and integrality specializations have been proved, that further cases of the mixed-Hodge formula have been established, and that the full mixed-Hodge statement for all \(\mathcal M_\mu\) remains unproved. This is supplied dated curation derived from the OWR material, not an independently verified literature assessment. The source requests independent verification of time-sensitive status and cited claims.
\item The registry labels the statement \texttt{corrected\_verified} and records that it was checked against the official Oberwolfach report. This is a supplied source-verification claim; neither that check nor the cited sources have been independently verified for this draft.
\end{itemize}
\end{document}
